% Generated draft: requires mathematical and source review.
\documentclass[11pt]{article}
\usepackage[margin=25mm]{geometry}
\usepackage{fontspec}
\setmainfont{Latin Modern Roman}
\usepackage{amsmath,amssymb,mathtools,unicode-math}
\setmathfont{Latin Modern Math}
\usepackage{enumitem,needspace,xurl,hyperref}
\setlength{\parindent}{0pt}
\setlength{\parskip}{5pt}
\setlength{\emergencystretch}{3em}
\newcommand{\N}{\mathbb{N}}
\newcommand{\Z}{\mathbb{Z}}
\newcommand{\Q}{\mathbb{Q}}
\newcommand{\R}{\mathbb{R}}
\newcommand{\C}{\mathbb{C}}
\newcommand{\F}{\mathbb{F}}
\newcommand{\A}{\mathbb{A}}
\newcommand{\PP}{\mathbb{P}}
\newcommand{\E}{\mathbb{E}}
\newcommand{\eps}{\varepsilon}
\newcommand{\dd}{\,\mathrm{d}}
\begin{document}
\section*{UnsolvedMath ID: 701028}
\textbf{Source problem code:} catalogue entry
\subsection{Definitions and mathematical statement}
Let \(\Z_{\ge 0}=\{0,1,2,\ldots\}\). For subgroups \(H,K\) of a group, let \([H,K]\) denote the subgroup generated by the commutators \(h^{-1}k^{-1}hk\), where \(h\in H\) and \(k\in K\). Define the lower central series by \(\gamma_1(H)=H\) and \(\gamma_{i+1}(H)=[\gamma_i(H),H]\). The nilpotency class of a nilpotent group \(H\) is the least \(c\in\Z_{\ge0}\) such that \(\gamma_{c+1}(H)=1\); thus the trivial group has class zero. Define \(G^{(0)}=G\) and \(G^{(i+1)}=[G^{(i)},G^{(i)}]\). A group is soluble if \(G^{(d)}=1\) for some \(d\in\Z_{\ge0}\), and its derived length \(\mathrm{dl}(G)\) is the least such \(d\).

Consider every finite group \(G\) with nilpotent subgroups \(A,B\le G\) of classes \(\alpha,\beta\), respectively, satisfying the setwise factorization
\[
G=AB=\{ab:a\in A,\ b\in B\}.
\]
Neither subgroup is assumed normal. The supplied source attributes to O. H. Kegel and H. Wielandt (1961) the result that every such \(G\) is soluble. It also states that \(\mathrm{dl}(G)\le\alpha+\beta\) need not hold, referring to Archive problem 5.17.

\textbf{Question.} Does there exist a function \(f:\Z_{\ge0}^{\,2}\to\Z_{\ge0}\), independent of \(G,A,B\) and their orders, such that every factorization above satisfies
\[
\mathrm{dl}(G)\le f(\alpha,\beta)?
\]
The source asks for a ``(linear) function.'' The stronger linear-bound formulation is whether there exist universal constants \(C_0,C_1,C_2\in\R_{\ge0}\) such that every such factorization satisfies
\[
\mathrm{dl}(G)\le C_0+C_1\alpha+C_2\beta.
\]
Here the displayed affine bound makes explicit an interpretation of the parenthetical word ``linear''; the supplied wording does not specify whether an arbitrary bound or specifically a linear bound is the primary target.
\subsection{Short English statement}
For a finite group that is the product of two nilpotent subgroups, ask whether its derived length is bounded solely by the nilpotency classes of those subgroups, and whether a universal linear bound exists. The source states that the group is soluble but that the sum of the two classes need not bound its derived length.
\subsection{Sources}
\begin{itemize}
\item Supplied local UnsolvedMath registry record: permanent ID \(701028\), exact source problem code \texttt{catalogue entry}, title \emph{Kourovka Notebook Problem 14.43}. The registry labels the problem \texttt{open}; this is supplied status metadata, not an independently verified assessment.
\item \emph{The Kourovka Notebook (Unsolved Problems in Group Theory)}, problem 14.43, proposed by L. S. Kazarin in the 14th issue (1999). The supplied record identifies the 21st issue edition (2026), \texttt{arXiv:1401.0300v48}. The collection metadata names E. I. Khukhro and V. D. Mazurov as editors.
\item The source attributes solubility of finite products of two nilpotent subgroups to O. H. Kegel and H. Wielandt (1961), without supplying individual publication details.
\item The source refers to Archive problem 5.17 for failure of the bound \(\mathrm{dl}(G)\le\alpha+\beta\); the archive entry and its supporting references are not included in the supplied record.
\end{itemize}
\end{document}
